\subsection{P036 --- A semantic model-based fault detection approach for building energy systems}
\label{paper:P036}
\textbf{Authors:} Tingting Li, Yang Zhao, Chaobo Zhang, Kai Zhou, Xuejun Zhang\par
\textbf{Major Category:} Fault Detection, Diagnostics and Maintenance\par
\textbf{Secondary Category:} BIM and Semantic Information\par
\textbf{Keywords:} Semantic model, Fault detection, Ontology, Knowledge graph, Semantic rules, Building energy systems, Inference\par
\paragraph{主な主張}
building energy systemのdomain ontology, abstract semantic rules, target-system knowledge graphを組み合わせ, operation, control, equipment, sensorに関するfault detection solutionをsystem configurationに応じて生成し, inference chainとして診断根拠を提示するapproachを示した. \cite{Li:2022SemanticFDD}
\paragraph{新規性}
equipment-specificなsemantic FDD modelに対し, reusableなabstract ruleとgeneric skeletal ontologyを用いて, 異なるbuilding energy systemsへcustomized fault detection logicを展開できる構成を提示した点に新規性がある.
\paragraph{位置付け}
ontology, knowledge graph, reusable rulesによるsemantic reasoningをFDDへ適用する代表例であり, semantic modelをrule executionとdiagnostic inferenceへ接続する研究として位置付ける.
